\section{Anomaly constraints as theorem-level filters (gauge and mixed gravitational)}
\label{app:anomaly_theorem_filters}

This appendix upgrades the anomaly constraints used throughout the Standard-Model interface to a theorem-facing \emph{filter} on admissible EFT representatives.
The connection to the Ward/BRST layer is as follows:
nontrivial gauge anomalies are genuine obstructions to a BRST-invariant quantization map and therefore to Ward/Slavnov--Taylor identities (Appendix~\ref{app:ward_brst_consistency}).
Accordingly, anomaly cancellation is treated here as a non-negotiable, checkable constraint on the EFT field content registry.

\subsection{Local anomaly coefficients (4D chiral gauge theory)}

\paragraph{Field-content input.}
Fix a gauge group $G$ and a finite set of left-handed Weyl fermions in representations of $G$.
In this paper, the relevant registry is the closed $18$-multiplet set (three generations) together with the gauge group $U(1)_Y\times SU(2)_L\times SU(3)_c$ (Section~\ref{sec:sm_labeling_closure} and Proposition~\ref{prop:channel_to_gauge}).
Scalars do \emph{not} contribute to gauge anomalies in 4D, so the Higgs sector does not affect the local anomaly cancellation conditions (it matters for EFT feasibility and beta functions but not for chiral gauge anomalies).

\paragraph{Minimal anomaly list used as filters.}
For $G=U(1)_Y\times SU(2)\times SU(3)$, the standard local anomaly constraints can be expressed by the vanishing of the following coefficients (in a left-handed basis; sums include multiplicities from color and weak isospin components):
\begin{itemize}
\item \textbf{mixed gravitational--hypercharge anomaly:} $\sum Y$,
\item \textbf{cubic hypercharge anomaly:} $\sum Y^3$,
\item \textbf{mixed non-abelian anomalies:} $SU(2)^2U(1)$ and $SU(3)^2U(1)$ coefficients,
\item \textbf{pure non-abelian anomalies:} $SU(3)^3$ (and $SU(2)^3$, which vanishes in the SM matter representations used here).
\end{itemize}

\begin{definition}[Anomaly coefficients (filter form)]\label{def:anomaly_coefficients_filter_form}
\MathTag In a left-handed Weyl basis, define the following anomaly coefficients for the $U(1)_Y\times SU(2)\times SU(3)$ factors:
\begin{align*}
A_{\mathrm{grav}\text{-}Y} &:= \sum_{\text{LH Weyl }f} Y_f,\\
A_{Y^3} &:= \sum_{\text{LH Weyl }f} Y_f^3,\\
A_{2^2Y} &:= \sum_{\text{LH Weyl }f} Y_f\,T_2(R_f),\\
A_{3^2Y} &:= \sum_{\text{LH Weyl }f} Y_f\,T_3(R_f),\\
A_{3^3} &:= \sum_{\text{LH Weyl }f} \sigma_f\,T_3(R_f),
\end{align*}
where $T_i(R)$ is the Dynkin index for the representation $R$ of $SU(i)$ (with the fundamental normalized as $T(\text{fund})=1/2$), and $\sigma_f=+1$ for color triplets and $\sigma_f=-1$ for color anti-triplets when expressing the $SU(3)^3$ anomaly in a left-handed basis.
\end{definition}

\begin{proposition}[Local anomaly cancellation filter]\label{prop:local_anomaly_filter}
\MathTag If any of the coefficients in Definition~\ref{def:anomaly_coefficients_filter_form} is nonzero, then the corresponding chiral gauge theory is anomalous and cannot admit a BRST-invariant quantization that preserves the Ward/Slavnov--Taylor identities for gauge symmetry.
Equivalently, a necessary condition for Assumption~\ref{ass:brst_invariant_quantization} is that all listed coefficients vanish.
\end{proposition}
\begin{proof}
This is the standard anomaly obstruction statement in 4D gauge theory:
gauge anomalies correspond to nontrivial BRST cohomology classes (Appendix~\ref{app:ward_brst_consistency}, Remark~\ref{rem:trivial_nontrivial_anomalies}) and obstruct BRST-invariant quantization.
The coefficients above encode the relevant local anomaly polynomials for the product group factors.
\end{proof}

\subsection{Application to the closed labeling content}

The paper already records explicit anomaly sums and their vanishing for the Standard Model hypercharge assignments (Section~\ref{sec:sm_labeling_closure}).
We restate the conclusion here as a theorem-level filter outcome for the fixed field-content registry.

\begin{theorem}[Anomaly-free status of the closed labeling fermion content]\label{thm:closed_labeling_anomaly_free}
\MathTag Under the PDG convention $Q=T_3+Y$, the chiral fermion content used by the closed labeling interface (three generations of $Q_L,u_R,d_R,L_L,e_R$ and a sterile singlet $\nu_R$ with $Y=0$) satisfies:
\[
A_{\mathrm{grav}\text{-}Y}=A_{Y^3}=A_{2^2Y}=A_{3^2Y}=A_{3^3}=0.
\]
In particular, adding a sterile $\nu_R$ with $Y=0$ does not change any anomaly sums.
\end{theorem}
\begin{proof}
The mixed gravitational and cubic hypercharge sums, as well as the mixed non-abelian anomaly sums, are computed explicitly for one generation in Lemma~\ref{lem:sm_anomaly_sums}, and vanish.
The statement that adding $\nu_R$ with $Y=0$ preserves these cancellations is Proposition~\ref{prop:anomaly_nur}.
The $SU(3)^3$ cancellation is the standard vectorlike cancellation between triplet and anti-triplet content per generation in a left-handed basis, and is included in the standard anomaly-cancellation package referenced there.
\end{proof}

\subsection{Global $SU(2)$ anomaly (Witten parity constraint)}

In addition to local perturbative anomalies, $SU(2)$ gauge theories in four dimensions have a global consistency condition (Witten anomaly) that depends on the parity of half-integer isospin representations.
The present manuscript uses this condition as a finite, non-negotiable filter on admissible extensions involving additional $SU(2)$ doublets.

\begin{proposition}[Global $SU(2)$ consistency filter (parity of doublets)]\label{prop:su2_global_anomaly_filter}
\MathTag In a four-dimensional $SU(2)$ gauge theory, a necessary condition for global consistency is that the total number of left-handed Weyl fermion doublets (more precisely: the number of half-integer isospin representations counted with multiplicity) is even.
\end{proposition}
\begin{proof}
Standard (Witten global $SU(2)$ anomaly); see \cite{Witten1982SU2Anomaly}.
\end{proof}

\begin{remark}[Why this matters for the minimal $\nu_R$ extension]\label{rem:witten_filter_nur}
\AuditTag The minimal extension used in this paper adds $\nu_R$ as a sterile singlet with $Y=0$.
It does not add any new $SU(2)$ doublets and therefore does not affect the global $SU(2)$ parity constraint.
By contrast, any extension that adds additional $SU(2)$ doublets must be audited against Proposition~\ref{prop:su2_global_anomaly_filter}.
\end{remark}

\subsection{Connection back to Ward/BRST theoremization}

\AuditTag The role of this appendix is to turn anomaly constraints into explicit \emph{filters} on admissible EFT registries:
if a proposed field-content registry violates the local anomaly filter (Proposition~\ref{prop:local_anomaly_filter}) or the global $SU(2)$ parity filter (Proposition~\ref{prop:su2_global_anomaly_filter}), then Ward/BRST consistency at the quantum level cannot be imposed without modifying the registry (i.e.\ adding compensating matter) or changing the scope.
This is the precise sense in which anomaly cancellation is a theorem-facing prerequisite for T1-2 and for any later scattering/observable equivalence claims.

